\section{Analysis}
\label{sec:analysis}

The guarantee of Eq.~\eqref{eq:crc} is the risk-control result of
\citet{angelopoulos2023conformal}; we state it in our notation because the rest
of the section builds on it, and we then derive the two consequences that
decide how the gate should be used in practice.

\begin{proposition}[Harm budget]
\label{prop:budget}
Let $z_1,\ldots,z_n$ be exchangeable decision states and let
$\ell_\lambda$ be the loss of Eq.~\eqref{eq:gate-loss}, bounded in $[0,1]$ and
monotone non-increasing in $\lambda$. For the threshold $\hat\lambda$ of
Eq.~\eqref{eq:crc}, $\E[R(\hat\lambda)]\le\alpha$, where the expectation is
over the calibration draw.
\end{proposition}

\begin{proof}
The corrected empirical risk
$\frac{n}{n+1}\widehat R_n(\lambda)+\frac1{n+1}$ upper-bounds
$\E[\widehat R_n(\lambda)]$ for every fixed $\lambda$ by
$\E[\widehat R_n(\lambda)]=\frac{n}{n+1}R(\lambda)$. Because the loss is
monotone, the feasible set $\{\lambda:\frac{n}{n+1}\widehat R_n(\lambda)+
\frac1{n+1}\le\alpha\}$ is an upper set, its infimum is attained, and the
standard monotone-loss argument gives $\E[R(\hat\lambda)]\le\alpha$.
\end{proof}

\begin{corollary}[Selective risk]
\label{cor:selective}
Under the conditions of Proposition~\ref{prop:budget}, the harmful share among
accepted decisions satisfies $\mathrm{SR}\le\alpha/C$ for any realised
coverage $C>0$.
\end{corollary}

\begin{proof}
$\mathrm{SR}=R/C$ by Eq.~\eqref{eq:selective-risk} and $R\le\alpha$ in
expectation.
\end{proof}

Corollary~\ref{cor:selective} says that a budget on harmful selections does not
by itself bound the harmful share of what remains: a rule that accepts few
decisions can have a high harmful share while keeping the mass small. Both
quantities therefore belong in a report, which is why the experiments below
report coverage and harmful share together.

\subsection{Why interval certification collapses}

The alternative to a budget is certification: build an interval around each
score and act only if the interval excludes harm. With a split-conformal
symmetric interval and nominal level $\alpha$, the radius is the
$(1-\alpha)$ quantile of the absolute calibration residual,
\begin{equation}
 q_\alpha=\operatorname{Quantile}_{1-\alpha}\{|s(z,j)-\marginal(j\mid A)|\},
 \label{eq:radius}
\end{equation}
and the gate accepts when $\max_j s(z,j)-q_\alpha>0$.

\begin{proposition}[Coverage ceiling of interval certification]
\label{prop:interval}
Let $M=\max_j s(z,j)$ and let $Q$ be the calibration radius. The accepted
fraction of the interval gate is $\Pr[M>q_\alpha]$, so if $q_\alpha\ge
\operatorname{ess\,sup}M$ the gate abstains everywhere.
\end{proposition}

\begin{proof}
The gate accepts exactly when $M>q_\alpha$.
\end{proof}

The proposition is trivial as stated, but it identifies the failure mode
precisely. The radius is a quantile of the estimation error, while the
decision compares it with the value being estimated. When the score is
inflated relative to the utility, the error is large in the same units and the
radius grows with it. In the deployment states we study, the standard
deviation of the score is about twice that of the true marginal utility and
the residual median is larger than the utility of a typical candidate, so the
radius at small $\alpha$ exceeds nearly every score and the gate abstains on
almost all decisions. Budget control does not have this failure mode, because
the threshold adapts to where the accepted mass is rather than to the width of
the error distribution.

\subsection{Transfer of a calibrated threshold}

A threshold calibrated on one distribution is used on another whenever the
deployment setting differs from the calibration setting. The loss of the gate
is a bounded function of the state, so a discrepancy between the two residual
distributions translates directly into a discrepancy in risk.

\begin{proposition}[Shift margin]
\label{prop:shift}
Suppose the deployment distribution $P$ and the calibration distribution $Q$
satisfy $\sup_\lambda|R_P(\lambda)-R_Q(\lambda)|\le\delta$. Then a threshold
$\hat\lambda$ calibrated on $Q$ with budget $\alpha$ satisfies
$R_P(\hat\lambda)\le\alpha+\delta$ whenever the calibration guarantee holds on
$Q$.
\end{proposition}

\begin{proof}
$R_P(\hat\lambda)\le R_Q(\hat\lambda)+\delta$ and
$R_Q(\hat\lambda)\le\alpha$ by Proposition~\ref{prop:budget}.
\end{proof}

The proposition makes the transfer question empirical: how large is $\delta$
when calibration and deployment differ by dataset or by training run? Section
\ref{sec:experiments} estimates it directly by calibrating on some datasets or
runs and evaluating on others. A method that controls risk but not the shift
margin would be of little practical use, so this is the quantity we report
alongside the budget.
